\section{Limits of simpler amplification schemes}\label{sec:barriers}

The majority construction addresses a geometric difficulty with tensor
products: a low-rank check matrix can create short paths between its
own coordinate and equation ports. The following finite results delimit
two simpler constructions and bound the exponent of the present one.

\begin{proposition}[Literal tensor checks]\label{prop:tensor-barrier}
\formal{TapeStorageRates/TensorObstruction.lean}{TapeStorageRates.active_cross_union_subset_shortPaired,TapeStorageRates.tensor_shortPaired_card_ge_nullities}
Let $A$ and $B$ be square matrices over any field, each with diagonal
one. Let $I_A$ and $O_A$ be the sets of rows and columns of $A$ having
a nonzero off-diagonal entry, and define $I_B,O_B$ likewise. For the
check matrix $H=A\otimes B$, use its displayed row-coordinate pairing
and helpers $N_H(v)=\{w\ne v:H_{vw}\ne0\}$.
Every coordinate in
\[
 (I_A\times O_B)\ \cup\ (O_A\times I_B)
\]
has paired distance at most three. Consequently at least
$(\dim\ker A)(\dim\ker B)$ coordinates have paired distance at
most three.
\end{proposition}
\begin{outline}
For $a\in I_A$ and $c\in O_B$, choose $b\ne a$ and $d\ne c$ with
$A_{ab}\ne0$ and $B_{dc}\ne0$. The diagonal-one entries give the walk
\[
 L_{(a,c)}\;--\;R_{(a,d)}\;--\;L_{(b,c)}\;--\;R_{(a,c)}.
\]
Its middle edge uses $A_{ab}B_{dc}\ne0$. Transposing both matrices
and swapping incidence ports proves the other crossed product.
The nonzero rows of $I-A$ are contained in $I_A$, so
$\dim\ker A\le\operatorname{rank}(I-A)\le|I_A|$ by rank
subadditivity. Transposition gives $\dim\ker A\le|O_A|$, and
similarly for $B$. Counting $I_A\times O_B$ yields the stated lower
bound.
\end{outline}

This obstruction concerns the literal Kronecker checks with their
displayed pairing. It makes no claim about a different check matrix,
a different pairing, or a nonlinear presentation of the same code.
The signed majority construction avoids this particular obstruction
by cancelling the even-change support terms.

\paragraph{Local projection deficits.}
Let $E$ and $V$ be finite sets of edge coordinates and vertices, with endpoint maps
$\ell,r:E\to V$ satisfying $\ell(e)\ne r(e)$. Distinct coordinates
may have the same endpoints. Put
$S_v=\{e:\ell(e)=v\text{ or }r(e)=v\}$ and $d_v=|S_v|$.

\begin{proposition}[A local-code and two-slot ceiling]\label{prop:tanner-barrier}
\formal{TapeStorageRates/TannerCeiling.lean}{TapeStorageRates.tanner_logb_card_le_of_local_sizes,TapeStorageRates.tanner_three_quarters_of_local_sizes_and_loads}
Let $C\subseteq\mathcal A^E$ be nonempty, with $|\mathcal A|=q\ge2$.
Suppose real numbers $t_v$ satisfy the actual projection bounds
\[
 |\pi_{S_v}(C)|\le q^{d_v-t_v}\qquad(v\in V),
\]
where $\pi_{S_v}$ restricts a word to $S_v$. Then
\begin{equation}\label{eq:tanner-deficit}
 \log_q|C|\le |E|-\tfrac12\sum_v t_v.
\end{equation}
If, in addition, there is an assignment
$o:E\to V\times\{0,1\}$ with
$|o^{-1}(v,i)|\le t_v$ for every $v,i$, then
$\log_q|C|\le3|E|/4$.
No linearity, regularity or girth assumption is required, and the
assignment need not send an edge to one of its endpoints.
\end{proposition}
\begin{outline}
For a uniform word $X$ of $C$, the stars cover every coordinate
exactly twice. The submodular-cover inequality gives
$2H(X)\le\sum_v H(X_{S_v})$.
Each summand is at most $(d_v-t_v)\log q$, and
$\sum_vd_v=2|E|$, proving \cref{eq:tanner-deficit}.
For the second assertion, every edge is assigned once, so
$|E|=\sum_{v,i}|o^{-1}(v,i)|\le2\sum_vt_v$. Substitute this
inequality into \cref{eq:tanner-deficit}. The Lean proof reuses the
uniform-code entropy and submodular-cover formalizations accompanying
\cite{PaperTwo}.
\end{outline}

For example, in a two-sheet construction, assigning each selected
lifted edge to its earlier endpoint is such a two-slot assignment;
\lean{TapeStorageRates.tanner_three_quarters_of_local_sizes} checks
that specialization. The ceiling applies when both the load bounds
and the projection bounds hold for the code actually used. Support
conditions on selected checks alone do not establish those projection
bounds for their full kernel. Thus this result rules out that local-code
route to rates approaching one, without claiming to rule out all
check-selection constructions.

\begin{proposition}[The amplification exponent is below one half]\label{prop:exponent}
\formal{TapeStorageRates/MajorityExponent.lean}{TapeStorageRates.matrix_card_le_support_mul_rank,TapeStorageRates.OddGirthLinearSeed.rate_le_degree_fraction,TapeStorageRates.Majority.decay_eq_neg_half_log,TapeStorageRates.Majority.exponent_mem_Ioo}
A matrix $H$ on $N$ coordinates, with nonzero diagonal and at most
$s$ nonzero entries per row, satisfies $N\le s\operatorname{rank}H$.
Every seed in \cref{def:seed} therefore satisfies
$\rho\le D/(D+1)$. For an above-half seed and the constants of
\cref{eq:construction-constants}, one has $0<\alpha<1/2$.
\end{proposition}
\begin{outline}
View $H$ as a linear map. Its restriction to each single-coordinate
input has rank one, because the corresponding diagonal entry is
nonzero. The scalar case of the shared block-rank theorem
\lean{LeanMachine.SparseProjection.sum_column_ranks_le} bounds the
sum of those ranks by $s\operatorname{rank}H$, giving the first
inequality.
Apply it to a seed check $I-A$, whose rows have support at most $D+1$.
Since $U\le\ker(I-A)$, rank--nullity gives
$\rho\le D/(D+1)$. Thus $(1-\rho)(D+1)\ge1$; together with
$4\rho>2$, this gives $4\rho(1-\rho)>1/(D+1)$. Taking logarithms
in \cref{eq:decay} gives $\lambda<\tfrac12\log(D+1)$, and
positivity yields $0<\alpha<1/2$.
\end{outline}

The exponent bound is consistent with the elementary inverse-degree
rate gap. It limits this particular amplification guarantee; it is
not an upper bound on the rate of arbitrary constructions. For the
explicit seed, $\lambda=-\tfrac12\log(1-(7/1613)^2)$ and $D$ is very
large. The resulting exponent is positive but small. Its importance
here is asymptotic: every fixed inverse power of $\Delta$ decays
faster than the tape theorem's iterated-logarithm allowance.
